% Note accompanying main.tex: the asymmetry statement and what each of its ingredients rests on. Written 2026-10-07.
% It proves nothing new: it assembles. Every status line below is to be kept true when an ingredient changes.
% Sources read for this note (read, not rebuilt): conway99-c3-experimental (README.md, c3formal/README.md,
% c3formal/Challenge.lean, pdf/preuve_c3.tex) and conway-C7 (README.md, lean/Challenge.lean, tex/order7.tex);
% CI runs listed with `gh run list` on 2026-10-07.
% Author: Dilan Djematene alone (his instruction of 07/10); AI assistance in the Acknowledgements, no model named.
% The links are printed (his instruction); the repositories are private since 07/10 and must be opened before circulation.
% TODO: the statements attributed to [MM04],
% [Mak10], [B09], [BL11] follow the literature section of the order-3 manuscript and were not re-read in the papers.
\documentclass[11pt]{amsart}
\usepackage[utf8]{inputenc}
\usepackage[T1]{fontenc}
\usepackage{lmodern}
\usepackage{amsmath,amssymb,amsthm}
\usepackage[margin=1.15in]{geometry}
\usepackage{booktabs}
\usepackage{enumitem}
\usepackage[hidelinks]{hyperref}

\newtheorem{theorem}{Theorem}
\newtheorem{corollary}[theorem]{Corollary}
\theoremstyle{remark}
\newtheorem{remark}[theorem]{Remark}
\newcommand{\Z}{\mathbb{Z}}
\newcommand{\G}{\Gamma}

\title[A Conway 99-graph is asymmetric]{A Conway $99$-graph, if it exists, is asymmetric:\\ statement and dependencies}
\author{Dilan Djematene}
\date{Draft of 7 October 2026}

\begin{document}

\begin{abstract}
A Conway $99$-graph is a strongly regular graph with parameters $(99,14,1,2)$; whether one exists is open. Assembling
published results with three recent manuscripts, which exclude automorphisms of order $2$, $3$ and $7$, one obtains that
such a graph has no nontrivial automorphism. This note states that consequence and lists, prime by prime, the result it
uses, the nature of its proof (by hand, by computer, formalised or not) and its status. It contains no new mathematics. The statement is conditional on these ingredients as they stand; in particular the
manuscript on order $2$ has one validation of its computations still to be completed.
\end{abstract}

\maketitle

\section{Statement}

Let $\G$ be a strongly regular graph with parameters $(99,14,1,2)$ and $G=\mathrm{Aut}(\G)$.

\begin{theorem}\label{thm:asym}
Assume the results \textup{(A)}--\textup{(E)} of Section~\ref{sec:ing}. Then $G$ is trivial: a strongly regular graph with
parameters $(99,14,1,2)$, if it exists, is asymmetric.
\end{theorem}

\begin{proof}
If $G\neq1$, it has an element of prime order $p$ (Cauchy). By (A), $p\in\{2,3,7,11\}$. These four primes are excluded by
(B), (C), (D) and (E) respectively.
\end{proof}

\section{The ingredients}\label{sec:ing}

\begin{enumerate}[label=\textup{(\Alph*)},leftmargin=2.4em]
\item \emph{$|G|$ divides $2\cdot3^3\cdot7\cdot11$.} Makhnev and Minakova \cite{MM04}, by character theory. Published;
no computer.
\item \emph{No automorphism of order $2$.} The paper \cite{C2}. The first half (an involution gives a frame of $51$
vectors in one of the $156$ odd unimodular lattices of rank $24$ without vectors of norm $1$) is proved by hand. The
second half (no such frame in any of the $156$ lattices) is computer-assisted, and verified as described in Section~9 of \cite{C2}; in the present draft of \cite{C2} one
validation, the recertification of a part of the transfers between lattices, remains to be completed. Not refereed; not
formalised.
\item \emph{No automorphism of order $3$.} The manuscript \cite{D3}: linear algebra over $\Z$ and the Eisenstein integers,
without computer. An automorphism of order $3$ has $0$ or $3$ fixed points (a count that reproves a result of \cite{MM04});
in both cases two lattices of eigenvectors must have determinants congruent modulo $3$, and they are not. Formalised
in Lean~4 with Mathlib (Appendix~\ref{app:lean}); the formalisation was read, not rebuilt, for this note.  Not refereed. Independently, the case of an automorphism of order $3$ with fixed points
is excluded by computer in \cite{B09,BL11}.
\item \emph{No automorphism of order $7$.} Behbahani and Lam \cite{B09,BL11}, by a computer enumeration of orbit matrices;
published. Without a computer the case was open: Cesarz and Woldar \cite{CW25} prove that $G\cong\Z_7$ if $7$ divides $|G|$
(published), which does not exclude an automorphism of order $7$. The manuscript \cite{D7} excludes it by hand, starting from the orbit structure and the counting identities of \cite{CW25}; it is formalised in Lean~4 with Mathlib
(Appendix~\ref{app:lean}; read, not rebuilt, for this note).  Not refereed.
\item \emph{No automorphism of order $11$.} Behbahani and Lam \cite{B09,BL11}, by computer; published. A report of the
MIT PRIMES programme \cite{LM14} gives an elementary argument, which was not verified for this note.
\end{enumerate}

\begin{table}[ht]
\centering\footnotesize
\begin{tabular}{lllll}
\toprule
prime & excluded by & method & formalised & status\\
\midrule
$p\ne2,3,7,11$ & \cite{MM04} & characters, by hand & no & published\\
$2$ & \cite{C2} & by hand, then computer-assisted & no & manuscript\\
$3$ & \cite{D3} & by hand & Lean 4 & manuscript\\
$7$ & \cite{BL11} & computer & no & published\\
    & \cite{D7}, using \cite{CW25} & by hand & Lean 4 & manuscript\\
$11$ & \cite{BL11} & computer & no & published\\
\bottomrule
\end{tabular}
\caption{What Theorem~\ref{thm:asym} uses.}
\end{table}

\section{Remarks}

\begin{remark}[What is known without any computer]
From \cite{MM04}, \cite{CW25}, \cite{D3} and \cite{D7}: $|G|$ divides $2\cdot11$, and if $|G|$ is even it divides $6$
\cite{CW25}; hence $|G|\in\{1,2,11\}$. The two remaining values are excluded by the computations of \cite{BL11} (order $11$)
and of \cite{C2} (order $2$).
\end{remark}

\begin{remark}[Before these manuscripts]
The published results \cite{MM04,Mak10,B09,BL11,CM20,CW25} leave $|G|\in\{1,2,3\}$, the open cases being an involution and a
fixed-point-free automorphism of order $3$. The manuscript \cite{D3} removes $3$ and \cite{C2} removes $2$; the results of
\cite{CM20} on the orders $6$ and $9$ are then no longer needed.
\end{remark}

\begin{remark}[Formalisation]
Theorem~\ref{thm:asym} itself is not formalised: \textup{(A)}, \textup{(B)} and \textup{(E)} are admitted. A formal
statement ``the automorphism group is trivial'' would need \textup{(B)} at least as a named hypothesis.
\end{remark}

\begin{remark}[Consequence for searches]
Every search for a Conway $99$-graph that prescribes a nontrivial automorphism is bound to fail: a proof of existence or of
nonexistence has to deal with a graph without symmetry.
\end{remark}

\appendix
\section{The formal statements}\label{app:lean}

\emph{Order $3$.} Formalised in
Lean~4 with Mathlib: the statement file \texttt{Challenge.lean} imports only Mathlib and contains
\begin{quote}\footnotesize
\texttt{theorem srg99\_no\_orderOf\_three \{V : Type*\} [Fintype V] (G : SimpleGraph V)}\\
\texttt{\ \ \ \ [DecidableRel G.Adj] (hG : G.IsSRGWith 99 14 1 2) ($\sigma$ : G $\simeq$g G) :}\\
\texttt{\ \ \ \ orderOf $\sigma$ $\neq$ 3}
\end{quote}
(namespace \texttt{Conway99}; Lean \texttt{v4.35.0-rc2}). According to the repository, the proofs use no \texttt{sorry} and
only the axioms \texttt{propext}, \texttt{Classical.choice}, \texttt{Quot.sound}, and \texttt{lake comparator} checks that
the proved statements are those of the statement file; its verification workflow (clean rebuild, comparator, and five
further checkers) reports success on 23 September 2026, at commit \texttt{874903ff}, and the three later commits change no
Lean file. For this note the statement file and the description were read; the formalisation was not rebuilt.

\emph{Order $7$.} \cite{D7} is formalised in Lean~4 with Mathlib
(\texttt{theorem srg99\_no\_orderOf\_seven}, same form as above with $7$; no \texttt{sorry}, the same three axioms,
\texttt{lake comparator}; verification workflow successful on 2 October 2026 at commit \texttt{8c6ce690}, no Lean file
changed since; read, not rebuilt, for this note).

\section*{Acknowledgements}
This work was carried out with the assistance of AI systems (large language models), which were used to develop the
arguments, to write and check the programs, and to draft and cross-check the texts. The author directed the work and takes
full responsibility for its content.

\begin{thebibliography}{99}
\bibitem[B09]{B09} M.~Behbahani, \emph{On strongly regular graphs}, Ph.D. thesis, Concordia University, Montr\'eal, 2009.
\bibitem[BL11]{BL11} M.~Behbahani and C.~Lam, \emph{Strongly regular graphs with non-trivial automorphisms}, Discrete Math.
\textbf{311} (2011), no.~2--3, 132--144.
\bibitem[C2]{C2} D.~Djematene, \emph{A strongly regular graph with parameters $(99,14,1,2)$
has no automorphism of order two}, manuscript and supplement, October 2026,
\url{https://github.com/amazing-source/strongly-regular-graphs-research/tree/main/papers/conway99_order2}; programs, records and certificates available from the author on request.
\bibitem[CM20]{CM20} D.~Crnkovi\'c and M.~Maksimovi\'c, \emph{Construction of strongly regular graphs having an automorphism
group of composite order}, Contrib. Discrete Math. \textbf{15} (2020), no.~1, 22--41.
\bibitem[CW25]{CW25} P.~G.~Cesarz and A.~J.~Woldar, \emph{On the automorphism group of a putative Conway 99-graph}, Algebr.
Comb. \textbf{8} (2025), no.~2, 379--398.
\bibitem[D3]{D3} D.~Djematene, \emph{Le graphe de Conway $\mathrm{srg}(99,14,1,2)$ n'a pas
d'automorphisme d'ordre $3$: une obstruction par les discriminants}, manuscript with a formalisation in Lean~4, September
2026, \url{https://github.com/amazing-source/strongly-regular-graphs-research/tree/main/papers/conway99_order3}; Lean formalisation available from the author on request.
\bibitem[D7]{D7} D.~Djematene, \emph{A computer-free proof that a Conway $99$-graph has no
automorphism of order seven}, manuscript with a formalisation in Lean~4, October 2026,
\url{https://github.com/amazing-source/strongly-regular-graphs-research/tree/main/papers/conway99_order7}; Lean formalisation available from the author on request.
\bibitem[LM14]{LM14} S.~Lou and M.~Murin, \emph{On the strongly regular graph of parameters $(99,14,1,2)$}, MIT PRIMES
report, 2014.
\bibitem[Mak10]{Mak10} A.~A.~Makhnev, \emph{On automorphisms of distance-regular graphs}, J. Math. Sci. \textbf{166} (2010),
no.~6, 733--742.
\bibitem[MM04]{MM04} A.~A.~Makhnev and I.~M.~Minakova, \emph{On automorphisms of strongly regular graphs with parameters
$\lambda=1$, $\mu=2$}, Diskret. Mat. \textbf{16} (2004), no.~1, 95--104; English transl.: Discrete Math. Appl. \textbf{14}
(2004), no.~2, 201--210.
\end{thebibliography}

\end{document}
